\documentclass[11pt]{article}
\usepackage[letterpaper,margin=1in]{geometry}
\usepackage{amsmath,amssymb,booktabs,enumitem}
\usepackage[hidelinks]{hyperref}
\setlength{\parskip}{0.5em}\setlength{\parindent}{0pt}
\setlist{itemsep=0.15em,topsep=0.2em}
\numberwithin{equation}{section}
\newcommand{\ice}{\mathrm{ice}}
\newcommand{\air}{\mathrm{air}}
\newcommand{\Rbar}{\bar R}
\newcommand{\rvs}{\rho_{vs}^{I}}
\newcommand{\pd}[2]{\frac{\partial #1}{\partial #2}}
\title{A time--temperature scaling relation for isothermal sintering by vapor transport}
\author{enceladus\_DSM \quad \texttt{studies/keff\_sintering/temperature\_scaling/}}
\date{2026-10-13}
\begin{document}
\maketitle

\section{Statement of the result}

Consider an aggregate of ice grains of characteristic radius $\Rbar$ in air,
held at a uniform temperature $T$ with no imposed gradient, evolving by
sublimation, vapor diffusion and desublimation as described by the
phase-field model of the manuscript. Let $Q$ be any quantity determined by
the geometry of the ice (specific surface area, a neck width, the effective
conductivity divided by that of ice). Then
\begin{equation}
\boxed{\;Q(t;\,T) = \widehat Q\!\left(\frac{t}{\tau(T)};\ \Lambda(T)\right),\qquad
\tau(T) = \frac{\Rbar^{2}\,\beta(T)}{d_0(T)},\qquad
\Lambda(T) = \frac{L^{*}(T)}{\Rbar}\;}
\label{eq:result}
\end{equation}
where $\widehat Q$ depends on the initial geometry but not otherwise on
temperature. The time scale $\tau$ contains the saturation vapor density and
changes by a factor of 29 between $-40$ and $-5\,^{\circ}$C; the group
$\Lambda$ contains the vapor diffusivity and changes by 20\%. To the extent
that this weak dependence can be ignored,
\begin{equation}
Q(t;\,T) = Q\!\left(\frac{t}{a_T};\ T_{\mathrm{ref}}\right),
\qquad
a_T = \frac{\tau(T)}{\tau(T_{\mathrm{ref}})}
    \simeq \sqrt{\frac{T}{T_{\mathrm{ref}}}}\;\frac{\rvs(T_{\mathrm{ref}})}{\rvs(T)} ,
\label{eq:shift}
\end{equation}
that is, a change of temperature is a change of the unit of time. Sections
2--7 derive this; Section~8 states it with its numbers; Section~9 gives the
conditions under which it holds.

Temperatures in formulas are absolute. Symbols follow the manuscript: $\phi$
for the phase field, $\chi$ for curvature, subscripts ``ice'' and ``air''.

\section{The model}

The unknowns are the phase fields $\phi_\ice$ and $\phi_\air = 1 - \phi_\ice$,
the temperature $T$ and the vapor density $\rho_v$. The governing equations
(manuscript Eq.~1) are
\begin{align}
\pd{\phi_\ice}{t} &= -3 M_{\mathrm{sub}}\left[\frac{1}{\varepsilon}\,\phi_\ice(1-\phi_\ice)(1-2\phi_\ice) - \varepsilon\nabla^{2}\phi_\ice\right]
 + \alpha_{\mathrm{sub}}\,\phi_\ice^{2}\phi_\air^{2}\,\frac{\rho_v - \rvs(T)}{\rho_\ice} , \label{eq:pf}\\[0.5ex]
\rho c_p\,\pd{T}{t} &= \nabla\cdot\left(K\nabla T\right) - \rho\,L_{\mathrm{sub}}\,\pd{\phi_\air}{t} , \label{eq:heat}\\[0.5ex]
\frac{1}{\xi_v}\,\pd{(\phi_\air\rho_v)}{t} &= \nabla\cdot\left(\phi_\air D_v\nabla\rho_v\right) + \rho_\ice\,\pd{\phi_\air}{t} . \label{eq:vap}
\end{align}
Equation~\eqref{eq:vap} is written as the solver implements it
(\texttt{src/assembly.c}): the factor $\xi_v = 10^{-3}$ multiplies the
diffusion and source terms together. With $\xi_v = 1$ it is manuscript
Eq.~(1c).

The kinetic parameters of Eq.~\eqref{eq:pf} are set from physical ones
(\texttt{src/material\_properties.c}, following Moure \& Fu, 2024):
\begin{equation}
\lambda_{\mathrm{sub}} = \frac{a_1\,\varepsilon\,\rho_\ice}{d_0\,\rvs},\qquad
\tau_{\mathrm{sub}} = \frac{\varepsilon^{2}\left(\beta_{\mathrm{sub}} + \Delta\beta\right)}{d_0},\qquad
M_{\mathrm{sub}} = \frac{\varepsilon}{3\,\tau_{\mathrm{sub}}},\qquad
\alpha_{\mathrm{sub}} = \frac{\lambda_{\mathrm{sub}}}{\tau_{\mathrm{sub}}} ,
\label{eq:pfparams}
\end{equation}
with $a_1 = 5$. The physical inputs are
\begin{align}
d_0 &= \frac{\gamma\,m}{\rho_\ice\,k_B T} && \text{capillary length,} \label{eq:d0}\\
\beta_{\mathrm{sub}} &= \frac{1}{\alpha_c}\,\frac{\rho_\ice}{\rvs(T)}\sqrt{\frac{2\pi m}{k_B T}} && \text{kinetic coefficient,} \label{eq:beta}\\
\Delta\beta &= a_1 a_2\,\varepsilon\left(\frac{1}{D_T} + \frac{1}{D_v(T)}\right)\frac{\rho_\ice}{\rvs(T)} && \text{thin-interface term, } a_2 = 0.1581, \label{eq:dbeta}\\
D_v &= D_{v0}\left(\frac{T}{T_0}\right)^{1.81} && \text{vapor diffusivity,} \label{eq:Dv}
\end{align}
and $\rvs(T)$, the saturation vapor density over flat ice (manuscript
Eqs.~2--3). The condensation coefficient $\alpha_c$ is a constant.

\textbf{Temperature enters through} $\rvs(T)$, $D_v(T)$, $d_0 \propto 1/T$
and the explicit $\sqrt{T}$ in $\beta_{\mathrm{sub}}$. Densities, heat
capacities and conductivities are constants in the model.

\section{Assumptions}

\begin{enumerate}[label=A\arabic*.]
\item \textbf{Isothermal setting.} The mean temperature is uniform and
  constant, and there is no imposed gradient. Write $T = T_m + \vartheta$
  with $\vartheta$ the departure caused by latent heat. (Its size is found in
  Section~6; it is not assumed to vanish.)
\item \textbf{Sharp-interface limit.} The interface parameter is small
  compared with the radii of curvature, $\varepsilon \ll \Rbar$, so the
  phase-field equation~\eqref{eq:pf} may be replaced by its sharp-interface
  limit.
\item \textbf{Constant $\alpha_c$.} The condensation coefficient does not
  depend on temperature or supersaturation.
\end{enumerate}
Three further simplifications are \emph{not} assumed. They are obtained in
Section~6 by showing that the corresponding dimensionless groups are small.

\section{The sharp-interface problem}

Let $\Gamma(t)$ be the ice surface, $\mathbf n$ its unit normal pointing from
the ice into the pore, and $v_n$ its velocity along $\mathbf n$, so that
$v_n > 0$ when ice grows.

\subsection{Vapor in the pore}
In the pore $\phi_\air = 1$ and $\partial\phi_\air/\partial t = 0$, and at
uniform temperature $D_v$ is uniform. Equation~\eqref{eq:vap} becomes
\begin{equation}
\frac{1}{\xi_v}\,\pd{\rho_v}{t} = D_v\,\nabla^{2}\rho_v \qquad\text{in the pore.}
\label{eq:s-diff}
\end{equation}
In the ice $\phi_\air = 0$ and Eq.~\eqref{eq:vap} is empty.

\subsection{Mass balance at the surface}
Take a coordinate $n$ along $\mathbf n$ with the surface at $n = s(t)$,
$\mathrm ds/\mathrm dt = v_n$. Across the surface $\phi_\air$ steps from 0 to
1, $\phi_\air = H(n - s(t))$ with $H$ the unit step, so
\[
\pd{\phi_\air}{t} = -v_n\,\delta(n - s),\qquad
\pd{(\phi_\air\rho_v)}{t} = -v_n\,\rho_v\,\delta(n - s) + \phi_\air\,\pd{\rho_v}{t} .
\]
Integrate Eq.~\eqref{eq:vap} over $n$ from $s - \eta$ to $s + \eta$ and let
$\eta \to 0$. The bounded term $\phi_\air\,\partial\rho_v/\partial t$
contributes nothing. The flux term gives $\phi_\air D_v\,
\partial\rho_v/\partial n$ evaluated on the pore side (it vanishes on the ice
side, where $\phi_\air = 0$). Hence
\[
-\frac{1}{\xi_v}\,v_n\,\rho_v = D_v\,\pd{\rho_v}{n}\bigg|_{\Gamma} - \rho_\ice\,v_n ,
\]
that is,
\begin{equation}
D_v\,\pd{\rho_v}{n}\bigg|_{\Gamma} = \left(\rho_\ice - \frac{\rho_v}{\xi_v}\right) v_n .
\label{eq:s-mass}
\end{equation}
This is conservation of water at the surface: ice that grows ($v_n > 0$) is
supplied by vapor diffusing toward it, so $\rho_v$ increases with distance
from a growing surface.

\subsection{Kinetics at the surface}
The matched-asymptotic analysis of Eq.~\eqref{eq:pf} for $\varepsilon \ll
\Rbar$ (Kaempfer \& Plapp, 2009; Moure \& Fu, 2024), with the parameters of
Eq.~\eqref{eq:pfparams}, gives the Gibbs--Thomson condition
\begin{equation}
\frac{\rho_v - \rvs(T_\Gamma)}{\rvs(T_\Gamma)} = d_0\,\chi + \beta\,v_n
\quad\text{on } \Gamma,
\qquad
\beta \equiv \beta_{\mathrm{sub}} + \Delta\beta = \beta_{\mathrm{sub}}\,(1 + \delta) ,
\label{eq:s-gt}
\end{equation}
where $\chi$ is the curvature of the surface (positive for a convex grain),
$T_\Gamma$ the local surface temperature, and
\begin{equation}
\delta \equiv \frac{\Delta\beta}{\beta_{\mathrm{sub}}}
 = a_1 a_2\,\varepsilon\,\alpha_c\left(\frac{1}{D_T} + \frac{1}{D_v}\right)\sqrt{\frac{k_B T}{2\pi m}} .
\label{eq:delta}
\end{equation}
This step is quoted, not derived here: it is the property the model is
constructed to have, and the cited papers give the analysis.
The factor $\rho_\ice/\rvs$ is common to $\beta_{\mathrm{sub}}$ and
$\Delta\beta$ and cancels in $\delta$.

\subsection{Heat}
Away from the surface Eq.~\eqref{eq:heat} is conduction in each phase,
\begin{equation}
\rho_j c_{p,j}\,\pd{\vartheta}{t} = k_j\,\nabla^{2}\vartheta,\qquad j \in \{\ice,\ \air\} .
\label{eq:s-heat}
\end{equation}
Integrating Eq.~\eqref{eq:heat} across the surface in the same way as in
Section~4.2, with $\partial\phi_\air/\partial t = -v_n\,\delta(n-s)$ and the density
of the phase that forms, $\rho_\ice$ (Eq.~\ref{eq:heat} weights this term with
the mixture density; the difference is immaterial here because the group
it feeds, $G_T$ below, is negligible),
\begin{equation}
k_\ice\,\pd{\vartheta}{n}\bigg|_{\ice} - k_\air\,\pd{\vartheta}{n}\bigg|_{\air} = \rho_\ice\,L_{\mathrm{sub}}\,v_n ,
\label{eq:s-stefanT}
\end{equation}
with $\vartheta$ continuous. Growth releases latent heat, which is conducted
away into both phases.

\subsection{Coupling of heat to the kinetics}
The surface temperature enters Eq.~\eqref{eq:s-gt} through $\rvs(T_\Gamma)$.
With $T_\Gamma = T_m + \vartheta$ and $|\vartheta| \ll T_m$,
\[
\rvs(T_m + \vartheta) = \rvs(T_m)\left[1 + s\,\vartheta + O(\vartheta^{2})\right],
\qquad
s \equiv \frac{\mathrm d\ln\rvs}{\mathrm dT}\bigg|_{T_m} \simeq \frac{L_{\mathrm{sub}}\,m}{k_B T_m^{2}} ,
\]
the last form being the Clausius--Clapeyron relation ($s = 0.096$~K$^{-1}$ at
$-20\,^{\circ}$C). Define the supersaturation relative to flat ice at the
mean temperature,
\begin{equation}
\sigma \equiv \frac{\rho_v - \rvs(T_m)}{\rvs(T_m)} .
\label{eq:sigma}
\end{equation}
Then, to first order in the small quantities $\sigma$ and $s\vartheta$,
\[
\frac{\rho_v - \rvs(T_\Gamma)}{\rvs(T_\Gamma)}
= \frac{1 + \sigma}{1 + s\vartheta} - 1 = \sigma - s\,\vartheta ,
\]
and Eq.~\eqref{eq:s-gt} becomes
\begin{equation}
\sigma = d_0\,\chi + \beta\,v_n + s\,\vartheta \qquad\text{on } \Gamma,
\label{eq:s-gt2}
\end{equation}
with $d_0$ and $\beta$ evaluated at $T_m$. From here on $\rvs$, $d_0$,
$\beta$, $D_v$ denote values at $T_m$, and $T$ is written for $T_m$.

\subsection{Summary of the dimensional problem}
Unknowns: $\rho_v$ (equivalently $\sigma$) in the pore, $\vartheta$
everywhere, and the surface $\Gamma(t)$. Equations:
\eqref{eq:s-diff}, \eqref{eq:s-mass}, \eqref{eq:s-gt2}, \eqref{eq:s-heat},
\eqref{eq:s-stefanT}, with periodic boundaries.

\section{Scales and dimensionless variables}

Choose
\begin{equation}
\begin{aligned}
\mathbf x &= \Rbar\,\mathbf x' , &
\chi &= \frac{\chi'}{\Rbar} , &
\sigma &= \frac{d_0}{\Rbar}\,u , \\
v_n &= V\,v' ,\quad V \equiv \frac{d_0}{\Rbar\,\beta} , &
t &= \tau\,t' ,\quad \tau \equiv \frac{\Rbar}{V} = \frac{\Rbar^{2}\beta}{d_0} , &
\vartheta &= \Theta\,\vartheta' ,\quad \Theta \equiv \frac{\rho_\ice L_{\mathrm{sub}} V \Rbar}{k_\ice} .
\end{aligned}
\label{eq:scales}
\end{equation}
The reasons:
\begin{itemize}
\item $\Rbar$ is the only length in the geometry.
\item $d_0/\Rbar$ is the supersaturation at which a grain of radius $\Rbar$
  is in equilibrium (Eq.~\ref{eq:s-gt2} with $v_n = 0$, $\vartheta = 0$,
  $\chi = 1/\Rbar$), so $u$ is of order one.
\item $V$ is the velocity that the kinetic term $\beta v_n$ needs for it to
  be as large as the capillary term $d_0/\Rbar$.
\item $\tau$ is the time for a surface moving at $V$ to travel $\Rbar$.
\item $\Theta$ is the temperature difference that conducts the latent heat
  released at velocity $V$ across a distance $\Rbar$ of ice.
\end{itemize}
Derivatives transform as
\[
\nabla = \frac{1}{\Rbar}\nabla' ,\qquad
\pd{}{n} = \frac{1}{\Rbar}\pd{}{n'} ,\qquad
\pd{}{t} = \frac{1}{\tau}\pd{}{t'} ,
\]
and from Eq.~\eqref{eq:sigma}, $\rho_v = \rvs\left(1 + \dfrac{d_0}{\Rbar}u\right)$,
so that $\nabla\rho_v = \rvs\,\dfrac{d_0}{\Rbar^{2}}\,\nabla' u$ and
$\partial\rho_v/\partial t = \rvs\,\dfrac{d_0}{\Rbar\,\tau}\,\partial u/\partial t'$.

\section{Nondimensionalization}

\subsection{Kinetics}
Substituting~\eqref{eq:scales} in Eq.~\eqref{eq:s-gt2}:
\[
\frac{d_0}{\Rbar}\,u = d_0\,\frac{\chi'}{\Rbar} + \beta\,\frac{d_0}{\Rbar\,\beta}\,v' + s\,\Theta\,\vartheta' .
\]
Dividing by $d_0/\Rbar$,
\begin{equation}
u = \chi' + v' + G_T\,\vartheta' ,
\qquad
G_T \equiv \frac{s\,\Theta\,\Rbar}{d_0}
 = \frac{s\,\rho_\ice L_{\mathrm{sub}}\,\Rbar}{k_\ice\,\beta} ,
\label{eq:nd-gt}
\end{equation}
where $\Theta\Rbar/d_0 = \rho_\ice L_{\mathrm{sub}} V \Rbar^{2}/(k_\ice d_0)
= \rho_\ice L_{\mathrm{sub}}\Rbar/(k_\ice\beta)$ by the definition of $V$.

\subsection{Mass balance}
The left side of Eq.~\eqref{eq:s-mass} is $D_v\,\rvs\,(d_0/\Rbar^{2})\,
\partial u/\partial n'$. On the right, $\rho_v = \rvs(1 + O(d_0/\Rbar))$,
so
\[
D_v\,\rvs\,\frac{d_0}{\Rbar^{2}}\,\pd{u}{n'} = \rho_\ice\left(1 - \frac{\rvs}{\xi_v\,\rho_\ice}\right) V\,v' .
\]
Solving for $v'$ and inserting $V = d_0/(\Rbar\beta)$,
\[
v' = \frac{D_v\,\rvs\,d_0}{\rho_\ice\,\Rbar^{2}}\cdot\frac{\Rbar\,\beta}{d_0}\cdot\frac{1}{1 - m_v}\,\pd{u}{n'}
   = \frac{D_v\,\beta\,\rvs}{\rho_\ice\,\Rbar}\cdot\frac{1}{1 - m_v}\,\pd{u}{n'} ,
\qquad
m_v \equiv \frac{\rvs}{\xi_v\,\rho_\ice} .
\]
Define the length
\begin{equation}
L^{*} \equiv \frac{D_v\,\beta\,\rvs}{\rho_\ice}
 = (1 + \delta)\,\frac{D_v(T)}{\alpha_c}\sqrt{\frac{2\pi m}{k_B T}} ,
\label{eq:Lstar}
\end{equation}
the second form following from Eqs.~\eqref{eq:beta} and~\eqref{eq:s-gt}:
$\beta\,\rvs/\rho_\ice = (1+\delta)\,\alpha_c^{-1}\sqrt{2\pi m/(k_B T)}$.
\textbf{The saturation vapor density cancels in $L^{*}$.} With $\Lambda
\equiv L^{*}/\Rbar$,
\begin{equation}
v' = \frac{\Lambda}{1 - m_v}\,\pd{u}{n'} \qquad\text{on } \Gamma .
\label{eq:nd-mass}
\end{equation}

\subsection{Vapor diffusion}
Substituting in Eq.~\eqref{eq:s-diff}, both sides carry $\rvs\,d_0/\Rbar$:
\[
\frac{1}{\xi_v}\,\rvs\,\frac{d_0}{\Rbar\,\tau}\,\pd{u}{t'} = D_v\,\rvs\,\frac{d_0}{\Rbar^{3}}\,\nabla'^{2}u
\quad\Longrightarrow\quad
\frac{P}{\xi_v}\,\pd{u}{t'} = \nabla'^{2}u ,
\qquad
P \equiv \frac{\Rbar^{2}}{D_v\,\tau} = \frac{d_0}{D_v\,\beta} .
\]
$P$ is the time for vapor to diffuse a distance $\Rbar$ divided by the time
for the geometry to change.

\subsection{Heat}
Equation~\eqref{eq:s-heat} becomes $(\Rbar^{2}/(D_{T,j}\,\tau))\,
\partial\vartheta'/\partial t' = \nabla'^{2}\vartheta'$ with $D_{T,j} =
k_j/(\rho_j c_{p,j})$, and Eq.~\eqref{eq:s-stefanT}, divided by
$k_\ice\Theta/\Rbar = \rho_\ice L_{\mathrm{sub}} V$,
\[
\pd{\vartheta'}{n'}\bigg|_{\ice} - \frac{k_\air}{k_\ice}\,\pd{\vartheta'}{n'}\bigg|_{\air} = v' .
\]
So $\vartheta'$ is of order one, and it enters the kinetics only through the
term $G_T\,\vartheta'$ of Eq.~\eqref{eq:nd-gt}.

\subsection{Sizes of the groups}
Evaluated for the aggregate runs ($\Rbar = 50~\mu$m, $\varepsilon = 1~\mu$m,
$\alpha_c = 1.020\times10^{-3}$) by \texttt{check\_groups.py}:

\begin{center}\small\setlength{\tabcolsep}{4pt}\begin{tabular}{llrrrrr}\toprule
Group & Meaning & $-40\,^{\circ}$C & $-30\,^{\circ}$C & $-20\,^{\circ}$C & $-10\,^{\circ}$C & $-5\,^{\circ}$C\\\midrule
$\Lambda = L^{*}/\Rbar$ & attachment against diffusion & 2.50 & 2.64 & 2.78 & 2.93 & 3.00\\
$\delta$ & thin-interface term & 0.0200 & 0.0200 & 0.0199 & 0.0199 & 0.0199\\
$G_T$ & latent heat on saturation & $1\times10^{-4}$ & $3\times10^{-4}$ & $7\times10^{-4}$ & $2\times10^{-3}$ & $2\times10^{-3}$\\
$m_v = \rvs/(\xi_v\rho_\ice)$ & vapor mass in the balance & $1\times10^{-4}$ & $3\times10^{-4}$ & $9\times10^{-4}$ & $2\times10^{-3}$ & $4\times10^{-3}$\\
$P/\xi_v$ & vapor storage & $1\times10^{-9}$ & $3\times10^{-9}$ & $7\times10^{-9}$ & $2\times10^{-8}$ & $2\times10^{-8}$\\
$\Rbar^{2}/(D_T\tau)$ & heat storage & \multicolumn{5}{c}{$10^{-12}$ to $10^{-10}$}\\\bottomrule
\end{tabular}\end{center}

$G_T$, $m_v$, $P/\xi_v$ and the heat-storage group are all below
$4\times10^{-3}$. Setting them to zero:
\begin{itemize}
\item $P/\xi_v \to 0$: the vapor field is quasi-steady, $\nabla'^{2}u = 0$.
\item $m_v \to 0$: the mass balance is $v' = \Lambda\,\partial u/\partial n'$.
\item $G_T \to 0$: the heat problem decouples from the kinetics, and the
  evolution is that of an isothermal system. (The scale $\Theta$ is
  $2\times10^{-8}$ to $6\times10^{-7}$~K.)
\end{itemize}
$\delta$ is not small compared with these, but it is the same at every
temperature to three digits, so it rescales $\beta$ by a constant and has
already been absorbed in $\tau$ and $L^{*}$.

\section{The reduced problem}

\begin{equation}
\boxed{\;
\begin{aligned}
\nabla'^{2}u &= 0 && \text{in the pore,}\\
u &= \chi' + v' && \text{on } \Gamma' ,\\
v' &= \Lambda\,\pd{u}{n'} && \text{on } \Gamma' ,
\end{aligned}
\;}
\label{eq:reduced}
\end{equation}
with periodic boundaries, the surface $\Gamma'$ moving at normal velocity
$v'$ in the time $t' = t/\tau$. Eliminating $v'$, the two surface conditions
are the single condition $u - \Lambda\,\partial u/\partial n' = \chi'$.

Problem~\eqref{eq:reduced} contains one parameter, $\Lambda$. Given the
initial geometry in units of $\Rbar$, its solution $\Gamma'(t')$ is a
function of $t'$ and $\Lambda$ only.

\textbf{Meaning of $\Lambda$.} A vapor molecule travelling a distance $\ell$
between two surfaces meets a diffusive resistance proportional to $\ell$ and
an attachment resistance proportional to $L^{*}$. For $\Lambda \to \infty$
the field $u$ is uniform, $v' = \bar u - \chi'$ with $\bar u$ fixed by
conservation of ice ($\oint v'\,\mathrm d\Gamma' = 0$), and the problem has
no parameter at all. For $\Lambda \to 0$, $u = \chi'$ on the surface and
diffusion alone limits the rate. The aggregates have $\Lambda \approx 2.5$
to 3: attachment is the larger resistance for transport between neighboring
surfaces, and diffusion is not negligible.

\section{The time--temperature scaling relation}

\subsection{Statement}
Because Problem~\eqref{eq:reduced} determines $\Gamma'$ as a function of
$t/\tau$ and $\Lambda$, any geometric quantity $Q$, made dimensionless with
$\Rbar$, satisfies Eq.~\eqref{eq:result}:
\[
Q(t;\,T) = \widehat Q\!\left(\frac{t}{\tau(T)};\ \Lambda(T)\right).
\]
This applies to the specific surface area (as $\mathrm{SSA}\cdot\Rbar$), to
neck widths, and to $k_{\mathrm{eff}}/k_\ice$, which for fixed
$k_\air/k_\ice$ is a functional of the geometry alone.

\subsection{The time scale}
From Eqs.~\eqref{eq:scales}, \eqref{eq:d0}, \eqref{eq:beta}
and~\eqref{eq:s-gt},
\[
\tau = \frac{\Rbar^{2}\beta}{d_0}
= \Rbar^{2}\cdot\frac{1+\delta}{\alpha_c}\,\frac{\rho_\ice}{\rvs}\sqrt{\frac{2\pi m}{k_B T}}\cdot\frac{\rho_\ice\,k_B T}{\gamma\,m} ,
\]
and since $\dfrac{k_B T}{m}\sqrt{\dfrac{2\pi m}{k_B T}} = \sqrt{\dfrac{2\pi k_B T}{m}}$,
\begin{equation}
\boxed{\;\tau(T) = (1+\delta)\,\frac{\Rbar^{2}\,\rho_\ice^{2}}{\alpha_c\,\gamma}\;\frac{1}{\rvs(T)}\sqrt{\frac{2\pi k_B T}{m}}\;}
\label{eq:tau}
\end{equation}
The temperature dependence is $\tau \propto \sqrt{T}/\rvs(T)$.

\subsection{The shift factor}
Relative to a reference temperature,
\begin{equation}
a_T \equiv \frac{\tau(T)}{\tau(T_{\mathrm{ref}})}
 = \sqrt{\frac{T}{T_{\mathrm{ref}}}}\;\frac{\rvs(T_{\mathrm{ref}})}{\rvs(T)} ,
\label{eq:aT}
\end{equation}
the factor $(1+\delta)$ cancelling to $10^{-4}$. It does not depend on
$\Rbar$, $\alpha_c$, $\gamma$ or $\varepsilon$.

\begin{center}\begin{tabular}{lrrrrr}\toprule
 & $-40\,^{\circ}$C & $-30\,^{\circ}$C & $-20\,^{\circ}$C & $-10\,^{\circ}$C & $-5\,^{\circ}$C\\\midrule
$\rvs$ [$10^{-4}$ kg\,m$^{-3}$] & 1.055 & 3.123 & 8.487 & 21.39 & 33.11\\
$V$ [$\mu$m\,d$^{-1}$] & 0.029 & 0.083 & 0.221 & 0.546 & 0.838\\
$\tau$ [d] & 1747 & 603 & 226 & 91.5 & 59.7\\
$a_T$ ($T_{\mathrm{ref}} = -20\,^{\circ}$C) & 7.72 & 2.66 & 1 & 0.404 & 0.264\\
$L^{*}$ [$\mu$m] & 125.0 & 132.0 & 139.2 & 146.4 & 150.1\\
$\Lambda$ & 2.50 & 2.64 & 2.78 & 2.93 & 3.00\\\bottomrule
\end{tabular}\end{center}

Between $-40$ and $-5\,^{\circ}$C, $\tau$ changes by a factor of 29.3: a
factor of 31.4 from $\rvs$ and 0.93 from $\sqrt T$.

\subsection{Activation energy}
In the model $\rvs$ is proportional to the saturation vapor pressure (for
$P_{vs} \ll P_{\mathrm{atm}}$), which obeys the Clausius--Clapeyron relation,
so $\rvs \propto \exp(-L_{\mathrm{sub}} M_w/(R_g T))$ with $M_w$ the molar
mass and $R_g$ the gas constant. Then, from Eq.~\eqref{eq:tau},
\begin{equation}
E_a \equiv R_g\,\frac{\mathrm d\ln\tau}{\mathrm d(1/T)} = L_{\mathrm{sub}}\,M_w - \tfrac12 R_g T \approx 50~\mathrm{kJ\,mol^{-1}} .
\label{eq:Ea}
\end{equation}
The formula gives 49.9 at $-20\,^{\circ}$C; a fit of $\ln\tau$ against $1/T$
over the five temperatures gives 50.1.

\subsection{The residual dependence}
From Eq.~\eqref{eq:Lstar} and~\eqref{eq:Dv}, $L^{*} \propto D_v/\sqrt T
\propto T^{1.31}$, so $\Lambda$ changes by $(268.15/233.15)^{1.31} = 1.20$
across the range. This is the only way the vapor diffusivity enters. If the
solution of Problem~\eqref{eq:reduced} is insensitive to a 20\% change in
$\Lambda$ about $\Lambda \approx 2.8$, Eq.~\eqref{eq:result} reduces to
Eq.~\eqref{eq:shift}: curves at different temperatures coincide when plotted
against $t/a_T$. In the limit $\Lambda \to \infty$ the coincidence is exact.

\subsection{A consequence that does not involve time}
Eliminating $t/\tau$ between two geometric quantities $Q_1$ and $Q_2$ of the
same run gives a relation $Q_1 = F(Q_2;\,\Lambda)$ that does not contain
$\tau$ at all. Temperature enters it only through $\Lambda$. This is why
$k_{\mathrm{eff}}$ plotted against SSA for one packing is the same at every
temperature, and why that comparison can be tighter than the comparison of
rates.

\subsection{Relation to the clock used in the manuscript}
From Eq.~\eqref{eq:pfparams},
\begin{equation}
\tau_{\mathrm{sub}} = \frac{\varepsilon^{2}\beta}{d_0} = \left(\frac{\varepsilon}{\Rbar}\right)^{2}\tau ,
\qquad
\theta \equiv \frac{t}{\tau_{\mathrm{sub}}} = \left(\frac{\Rbar}{\varepsilon}\right)^{2}\frac{t}{\tau} .
\end{equation}
With $\varepsilon/\Rbar = 1/50$ in every aggregate run, $\theta = 2500\,
t/\tau$, so collapse against $\theta$ and collapse against $t/\tau$ are the
same statement, and $\tau_{\mathrm{sub}}(T)/\tau_{\mathrm{sub}}(T_{\mathrm{ref}})
= a_T$. Check: $226.3~\mathrm{d}/2500 = 7822$~s, the value of
$\tau_{\mathrm{sub}}$ at $-20\,^{\circ}$C in manuscript Table~2; the other
four temperatures agree to 0.03\%.

\section{Conditions of validity}

The relation rests on A1--A3 and on the smallness of the groups of
Section~6.5. It fails, or needs amending, when:
\begin{itemize}
\item \textbf{$\alpha_c$ depends on temperature.} Then $\tau$ acquires that
  dependence and so does $\Lambda$ (Eq.~\ref{eq:Lstar}); the first changes
  $a_T$, the second breaks the collapse if $\alpha_c$ changes much.
\item \textbf{A second transport mechanism acts}, such as surface
  diffusion. It adds a term to the surface velocity with its own time scale
  and temperature dependence; the ratio of the two time scales is a new
  group that varies with temperature.
\item \textbf{A temperature gradient is imposed.} $\rvs$ then varies in
  space at order $s\,\Rbar\,|\nabla T|\,\Rbar/d_0$ in units of $u$, which
  exceeds one for gradients above about $4$~K\,m$^{-1}$ at these sizes. The
  isothermal problem no longer applies.
\item \textbf{$\Lambda$ is of order one or less and varies more}: larger
  grains, a larger $\alpha_c$, or a wider temperature range. The diffusivity
  then matters in full. For the Molaro pairs ($\alpha_c = 0.1$, $\Rbar
  \approx 87~\mu$m) $\Lambda \approx 0.02$.
\item \textbf{The interface is not thin}, $\varepsilon/\Rbar$ not small
  compared with the local radius of curvature. The departure from the
  sharp-interface problem is the same at every temperature if
  $\varepsilon/\Rbar$ is fixed, so it changes $\widehat Q$ but not the
  collapse.
\end{itemize}

\section{Comparison with the simulations}

The 125 aggregate runs (five packings at each of five porosities, each at
five temperatures) test the relation with the full diffuse-interface model:
\begin{itemize}
\item Rates follow $\tau_{\mathrm{sub}}$ to within 4\%: a packing reaches a
  given SSA at times in the ratio $a_T$.
\item $k_{\mathrm{eff}}$ at matched SSA agrees across the five temperatures
  to within 0.27\% for every packing (median 0.07\%), the time-free
  consequence of Section~8.6.
\end{itemize}
A 20\% change in $\Lambda$ is of the right size to account for the 4\%; that
it is the cause has not been tested. The analysis shows that the collapse is
a property of the model for any starting geometry, and the simulations show
how closely the diffuse-interface implementation follows it.

\appendix
\section{Constants used}
$\rho_\ice = 919$~kg\,m$^{-3}$; $k_\ice = 2.29$, $k_\air =
0.02$~W\,m$^{-1}$\,K$^{-1}$; $L_{\mathrm{sub}} = 2.83\times10^{6}$~J\,kg$^{-1}$;
$\gamma = 0.109$~J\,m$^{-2}$; $m = 2.9915\times10^{-26}$~kg; $D_{v0} =
2.178\times10^{-5}$~m$^{2}$\,s$^{-1}$ at $T_0 = 273.15$~K; $D_T =
7.78\times10^{-6}$~m$^{2}$\,s$^{-1}$; $a_1 = 5$, $a_2 = 0.1581$; $\xi_v =
10^{-3}$. Every number in this note is produced by \texttt{check\_groups.py},
which also asserts that the closed form~\eqref{eq:tau} equals
$\Rbar^{2}\beta/d_0$ and that $(\varepsilon/\Rbar)^{2}\tau$ reproduces
manuscript Table~2.

\end{document}
